\todo{Lecture 12}
\begin{example}[Helium]
Soient $\gamma=\frac{5}{3}$, $T=293~\mathrm{K}$ et $m=4.16\cdot 10^{-27}~\mathrm{kg}$ alors $c_{\mathrm{He}}= 1004~\mathrm{ms}^{-1}$.
\end{example}

\subsection{Tuyaux d'orgues/flutes}

poly p.14

\section{Quelques mots sur les ondes lineaires a la surface d'un fluide parfait}

Les vagues a la surface de l'eau sont une combinaison d'onde tranverse et longitudinale. Si on neglige la viscosite et tension superficielle, en regime lineaire, on trouver la relation de dispersion suivante : 
$$
\omega^{2}=gk\tanh(kh)
$$
donc $\frac{\omega}{k}$ n'est pas constante (c-a-d onde avec dispersion).
\begin{description}
\item[Cas 1 (eau profonde avec $kh=\frac{2\pi h}{\lambda} \gg 1$)] Come $hk \gg 1$ alors $\tanh(kh)\approx 1$ et $\frac{\omega}{k}=\sqrt{ \frac{g}{k} }=\sqrt{ \frac{g}{2\pi} \lambda}$. Les ondes sur la mer (la houle), generee par le vent $\lambda<100$.

\item[Cas 2 (eau peu profonde $kh \ll 1$)] Comme $hk\ll 1$ alors $\tanh(kh)\approx kh$. Donc $\frac{\omega}{k}=\sqrt{ gh }$. Les ondes a la plage, ondes de tsunami.
\end{description}